\chapter{Conclusion}\label{ch:conclusion}

This thesis demonstrated the technical feasibility of a low-cost closed-loop sacrum/lower-back wearable system for monitoring pelvic-region movement during prolonged sitting and delivering detection-outcome-based haptic feedback in pilot seated contexts. The prototype integrated local IMU sensing, BLE data transmission, Android labelled field logging, PC-based researcher inspection, compact embedded detection logic, an inactivity timer, and local vibration feedback. Trend-level comparisons with a research-grade Shimmer IMU supported the XIAO-based prototype as a plausible sensing source for the intended binary movement-detection task, while the pilot field data showed that instructed pelvic rotation could be distinguished from static sitting and that a compact KNN model could run within the timing and memory constraints of the target wearable. However, the study provides feasibility evidence only: it does not establish clinical pelvic kinematics, spontaneous daily-use robustness, long-term wearing comfort, battery lifetime, or behavioural effectiveness of the haptic reminder.
